\subsection{A three-isogeny addition tensor and the Hermitian cubic}
\label{subsec:f0-theta-Hermitian-containment}

This isolated subsection is a candidate.  It is not part of the
delivered CP31 paper.  The purpose is to check an obstruction to using
a necessary Moore-determinant equation as a contradiction on actual
theta parameters.

\begin{proposition}[The symmetric two-rank-zero theta image]
\label{prop:f0-theta-Hermitian-containment}
\status{Candidate; complete proof supplied.}
Let \((A,M)\) be a principally polarized abelian surface of two-rank
zero over an algebraically closed field of characteristic two, with
\(M\) symmetric.  Put \(L=M^3\), and choose the symmetric
level-three theta basis \((e_h)_{h\in\mathbf F_3^2}\) used in the
typed Cartier--Moore identification.  Then
\[
                         \sum_h e_h(a)^3=0
                         \qquad(a\in A).
                                                        \tag{THC1}
\]
This is an identity of sections of \(L^3\), so the complete
cubic-theta image lies on the Hermitian cubic
\(\sum_h X_h^3=0\).
\end{proposition}
\begin{proof}
Fix \(0\ne\theta\in H^0(A,M)\), and consider
\[
 \eta:A\times A\longrightarrow A\times A,
       \qquad (z,a)\longmapsto(2z+a,z-a).
                                                        \tag{THC2}
\]
It is an etale isogeny with kernel the diagonal copy of \(A[3]\).
The symmetric theorem of the square and cube give the rigidified
isomorphism
\[
 \eta^*(M\boxtimes M^2)\simeq L^2\boxtimes L.
                                                        \tag{THC3}
\]
Indeed, the coefficients in the two variables are \(4+2=6\)
and \(1+2=3\), while the cross terms cancel.  Thus
\[
            S_3(z,a)=\theta(2z+a)\theta(z-a)^2
                                                        \tag{THC4}
\]
is a nonzero section of the right side of (THC3).

We first show that its first-variable section lies in the squared
copy of \(H^0(A,L)\).  Write \(F:A\to A^{(2)}\) and
\(V:A^{(2)}\to A\) for relative Frobenius and Verschiebung,
so \([2]=VF\).  On \(A^{(2)}\times A\), the line bundle
\[
 N_1=(Vp_1+p_2)^*M\ \otimes\ (p_1-Fp_2)^*M^{(2)}
                                                        \tag{THC5}
\]
has a section whose pullback by \(F\times1\) is exactly
(THC4).  Formula (THC3) identifies its pullback with that of
\(N_0=L^{(2)}\boxtimes L\).  The quotient
\(N_1N_0^{-1}\) is therefore in the kernel of pullback by the
proper torsor \(F\times1\).  Such a line is a character line
for \(\ker F\): a trivialization of its pullback has descent
functions in \(\mathcal O(\ker F)^*\), because the base product
is proper and connected, and the cocycle identity is exactly the
character identity.  Since \(A\) has two-rank zero,
\((\ker F)^D\) is connected and has only its identity geometric
point.  This character is trivial.  Hence \(N_1\simeq N_0\),
and the section in (THC4) belongs to
\[
   F^*H^0(A^{(2)},L^{(2)})\otimes H^0(A,L)
       =\langle e_h(z)^2\rangle_h\otimes H^0(A,L).
                                                        \tag{THC6}
\]
This uses the group-scheme character argument, not a basis of
geometric characters for the coordinate algebra of a connected
group scheme.

We next check the exact diagonal three-torsion action.  Pullback
through \(\eta\) gives a natural diagonal \(A[3]\)-linearization
on (THC3), under which \(S_3\) is fixed.  The finite theta group
of \(L\) supplies another such linearization: on a theta lift
\(R_t\), its operator is \(R_t^2\boxtimes R_t\).  Multiplying
a finite lift by a cube root of unity changes this operator by
its cube, hence does not change it; the cocycle scalars likewise
cancel.  These operators therefore define an actual action of
the diagonal \(A[3]\).

The ratio of the two linearizations is a character of \(A[3]\).
The rigidified isomorphism (THC3) respects simultaneous inversion
of \(z,a\), by symmetry and its value at \((0,0)\).  Consequently
this character has the same value at \(t\) and \(-t\).  Its
order divides both two and three, so it is trivial.  Thus the
section in (THC6) is fixed by the exactly specified diagonal
theta action.  In the symmetric Schr\"odinger basis, invariance
under the character Lagrangian forces a term
\(e_h(z)^2e_j(a)\) to have \(2h+j=0\), or \(j=h\).
Invariance under the translation Lagrangian makes all nine
coefficients equal.  Nonvanishing of \(S_3\) gives
\[
       \theta(2z+a)\theta(z-a)^2
                =c\sum_h e_h(z)^2e_h(a),\qquad c\ne0.
                                                        \tag{THC7}
\]
The equality is interpreted through the one rigidified line-bundle
isomorphism and thus has only one common scalar.

Finally \(\theta(0)=0\).  For an indecomposable polarization,
the previously proved theta/Jacobian dichotomy and two-rank-zero
genus-two normal form identify the symmetric theta with the curve
based at its unique Weierstrass point; it contains the origin.
There is no different symmetric representative up to isomorphism,
since the geometric two-torsion of \(\operatorname{Pic}^0(A)\)
is trivial.  For a product polarization both elliptic factors are
supersingular, and its symmetric theta is the union of the two
coordinate elliptic curves, also containing the origin.  Set
\(z=a\) in (THC7).  The left side is
\(\theta(3a)\theta(0)^2=0\), proving (THC1).
\end{proof}

\begin{remark}[Exact limitation of the calculation]
\label{rem:f0-theta-Hermitian-containment-scope}
\status{Candidate scope; no determinant-divisibility conclusion.}
In two-rank zero the actual parameter vector in the typed Moore
identification is \([c_h(a)]=[e_h(a)^2]\).  Thus (THC1), if
confirmed by independent audit, implies \(\sum_h c_h(a)^3=0\).
The necessary coefficient condition of that form in a Fermat
divisibility calculation would then be automatic on the actual
parameter surface.  The present argument does not say that the
Moore determinant is divisible by the Fermat cubic, and it does
not decide the remaining affine-line conditions.

For positive two-rank, the character line in (THC5)--(THC6) need
not be trivial.  This proof therefore makes no two-rank-one or
ordinary assertion.  All inputs used here are the previously
audited theorem-of-cube, theta/Jacobian, normal-form and
group-scheme descent arguments.  No new numerical computation
or external-source access is required or claimed.
\end{remark}
